Een monnik bood begaan
de jager voedsel aan
en wees waar hij moest gaan.

Hij kon zijn soep gaan eten
in een tijdversleten
predikstoel gezeten.

Emmy stond nog steeds krom.
Haar nekje in de kom
van de houten kolom

die je kon scharnieren.
Zo konden de klieren
hun kwelzucht botvieren

zonder enig verzet.
Het was wrokkige pret
verankerd in de wet.

De schandpaal had daarnaast
twee polsgaten ernaast.
Het publiek was verbaasd

dat haar kleine handen
aldaar niet belandden.
Met niets meer omhanden

lagen ze op haar rug,
die vormde een hangbrug
naar haar beentjes terug.

Een wasvrouw spitsvondig
deed haar schoonmaak grondig
en riep daarom mondig

“Benen wijd, vieze slet”.
In de schandpaal gezet
lukte dat echter net.

Emmy zo uitgewoond
werd al staande verschoond.
De wasvrouw werd beloond

met enkele muntjes.
Ze deed haar wasstuntjes
tot in de mikpuntjes.

Emmy reflecterend
en analyserend,
dacht toen concluderend

“Dit huwelijksritueel,
ik weet daarvan niet veel,
maar mist ze niet een deel,

namelijk mijn gezicht.”
Haar mond bleef echter dicht.
Het was ook niet haar plicht.

Ze wilde niet twisten
en geen tijd verkwisten
om als een braaf christen

in haar ere gekweld
snel te worden hersteld.
Haar zorg gerustgesteld

over seksueel onrein.
Een gevallen vrouw zijn
dat vond Emmy niet fijn.

Hij deed zijn huwelijksplicht
nadat zij was gezwicht.
Nu werd zij afgericht

conform deze sekte.
Het was geen perfecte
sfeer die men verstrekte,

eerder afstotelijk,
maar het was christelijk
en niet bespottelijk.

Als kennelijk alles
toch een illusie is
wat is er dan mis

met het fijne verschil
als een roze wensbril
laat zien wat Emmy wil?

Elke imperfectie
checkte de inspectie
strikt sectie voor sectie.

De viering vervolgde met de inspectie.
Dat deed men grondig sectie voor sectie.

Een oogarts bekeek haar bruine kijkers
Hij behoorde tot de azijnzeikers

toen hij een lichte afwijking constateerde
die niemand anders ietwat deerde.

Een tandarts inspecteerde haar witte tanden.
Overal voelde Emmy warme handen,

die checkten de vlekjes op haar witte huid.
Een dokter commandeerde; “steek je tong uit!”

“Ze heeft nog een vochtige rode lap”,
zei een naaister voor de grap.

In elk geintje
zit een seintje.

Dit uitlachen ging maar door.
Daar zorgde de hofnar wel voor.

Door wat er eerder was gebeurd
waren haar sproetjes besmeurd.

De nar veegde het op haar tong,
terwijl hij “smakelijk eten” zong.

Het besef kwam en de schok was groot,
dit is niet een bijzondere huwelijksboot.

Ze is niet meer dan gevangen wild,
waarmee je kunt doen wat je maar wilt.

De herder had haar goed bekeken
en haar met een lammetje vergeleken.

Die lam zou in de dienst worden geofferd.
Haar onschuld worden geslachtoffert.

Daarmee nam zij op zich alle zonden
Die uit de kerkviering ontstonden.

De barman goot wijn over het lam
het zette haar huid in vuur en vlam.

Haar hoeven waren voor de pedicure
die oordeelde; “mooi van nature!”.

Op haar kop kastanjebruine wol.
Een herder greep er een hand van vol.

Hij wilde het graag afscheren
met als excuus ongedierte te weren.

Een kapster wist hem in te dammen
en begon de klitjes eruit te kammen.

Een dierenarts controleerde op luis
en gaf aan dat zit wel pluis.

Daarop kneep een muzikant in haar tieten
om van haar ritmische kreuntjes te genieten

Van achteren greep een slager een borst
en legde zijn veel te kleine worst

in haar nog geboeide handjes
ze moest knijpen in de vetrandjes.

Een schilder tatoeerde op haar rechterdij
het gangbare teken voor vogelvrij

Een drukker brandde in haar linkerbil
de korte zin; “Ja Meester, ik wil!”

De smid leek op een dwerg
Zijn buik was als een berg.

Hij had een lange witte baard
Droeg aan zijn riem een zwaard.

Leek gemaakt uit steen en zand
foeilelijk met een beperkt verstand.

De smid wijkte zo af van het mensenras
dacht Emmy dat hij een dwerg was

Dat hij in het duister was verschenen
omdat zonlicht hem zou verstenen.

Emmy kende dwergen als zelfzuchtig
hij stal zowaar haar hanger hebzuchtig

De smid verving toch wat wijzer
haar halsband door een van ijzer.

Zorgde dat deze precies goed past
Smeedde er een ketting aan vast.

Een koene ridder van adel
gebruikte haar handjes als zadel

de wandelstok als beugel
en de ketting als teugel.

Zo wilde hij haar berijden.
Moesten haar beentjes strijden.

Ze bewaarde net hun evenwicht.
Bezweken later onder het gewicht.

Zo speelden alle stedelingen hun rol
en degradeerde Emmy tot snol.

“Zou deze komedie
deze tergende parodie

een goede afloop kennen?
Of zal ik moeten wennen

aan een verblijf in de hel.
Ik zie het allemaal wel!”

Of hoort dit bij een huwelijksritueel?
De priester verschijnt pas later

Daarbij verschoof haar aandacht
wat haar naar het heden bracht.

Ze was nu wel ouder dan achttien
en men had het op haar voorzien.

Tot nu toe had men haar kuisheid bewaard
haar drie gaatjes zorgvuldig gespaard.
